% !TEX root = ../main-es.tex
%TC:macro \limite [ignore,ignore]
%TC:macro \nota [ignore]
%TC:macro \TODO [ignore]
%TC:macro \parencite [ignore]
%TC:macro \textcite [ignore]
\section{Conclusiones}\label{sec:conclusiones}

Las cifras de este trabajo salen de 16.260 corridas hechas con la regla de costo corregida, y varias conclusiones de la versión anterior del modelo no sobrevivieron a la corrección. El cuadro que queda es más matizado que la frontera limpia que se conjeturaba al comenzar, y con 30 réplicas por combinación varias diferencias entre políticas caen dentro del error estándar, lo cual se indica en cada caso, del mismo modo que se indican las hipótesis que los datos rechazaron.

\begin{enumerate}[leftmargin=1.8em, itemsep=0.35em]
  \item \textbf{Con el mismo replanificador al mismo precio, reparar al ver la divergencia y esperar a que el plan falle casi empatan.} Con $r=3$, la diferencia entre \OL{} y \CL{} va de $-9$ a $+26$ puntos y en 13 de las 16 celdas dinámicas queda dentro de dos errores estándar, ningún factor barrido explica más del 10\,\% de su varianza y su signo cambia en la mitad de las celdas al sortear el orden de activación, de modo que monitorear el plan no compra casi nada sobre replanificar solo al fallar.
  \item \textbf{La frontera que sí aparece separa deliberar de no deliberar.} Cuando el mundo cambia rápido y deliberar es caro, el agente reactivo basado en modelo es la mejor política (98.3 a $\lambda=0.5$ contra 53.3 de la mejor que planifica con $c=2$, y 86.0 contra 56.0 a $\lambda=0.25$ con $c=5$), de forma que el costo de reescritura a partir del cual la ventaja se invierte existe, pero es el que hace dejar de convenir planificar (entre $c=0$ y $c=2$ a $\lambda=0.5$, y entre $c=2$ y $c=5$ a $\lambda=0.25$), y el orden de la taxonomía de agentes depende del entorno de tarea. Es la conclusión más estable del trabajo, pues sobrevive al orden de activación, aparece desde el tick 15 a $\lambda=0.5$ y crece con la abundancia de objetivos (a $\lambda=0.5$, de \dif{+13.0}{3.1} con 2 recursos a \dif{+96.0}{13.0} con 16).
  \item \textbf{Reparar en caliente vale claramente más que reiniciar en frío desde $\lambda=0.1$,} con $+16$ a $+22$ puntos a $\lambda=0.1$ y $0.25$ cuando $c$ está entre 2 y 5, $+32$ puntos ($+41$\,\%) a $\lambda=0.5$ con deliberación gratuita y hasta $+80.7$ con visión total, aunque la relación con $c$ no es monótona y la ventaja se anula donde ambas políticas colapsan. Se mantiene positiva con el orden sorteado, pero tarda en pagar (se separa del ruido desde el tick 105 a 130, cuando hay trabajo en curso que perder) y desaparece cuando los recursos escasean.
  \item \textbf{La observabilidad amplifica lo que decide el costo de deliberar.} Ver más sube el nivel de todas las políticas (de 47.0 a 261.0 a $\lambda=0$ entre $r=1$ y $r=\infty$), pero con $c=5$ la visión total no hace ganar al lazo cerrado (114.5 contra 129.5 del abierto, en promedio sobre $\lambda$), que solo domina con deliberación gratuita y oráculo perfecto.
  \item \textbf{Acotar las reparaciones proactivas evita el peor caso, pero no hay una cuota óptima que rescate al agente.} A $\lambda=0.5$, $B=3$ rinde 23.7 contra 10.0 del lazo cerrado, pero $B=0$ (el lazo abierto) rinde 27.0, y el presupuesto explica a lo sumo el 0.3\,\% de la varianza, por lo que un tope fijo y bajo de iteraciones no es la razón por la que el agente funciona.
  \item \textbf{La competencia del replanificador pesa más sobre quien replanifica menos, y en general pesa poco.} Subir $p$ de 0.5 a 1 lleva al lazo abierto de 36.7 a 61.0 a $\lambda=0.25$ y al cerrado solo de 48.3 a 59.7, pero $p$ explica a lo sumo el 1.5\,\% de la varianza de cada política y no deja huella en la telemetría de un motor en frío.
  \item \textbf{Las políticas se separan por ritmo, no por un evento.} Las entregas acumuladas y los ticks deliberando crecen casi en línea recta desde el tick 50, con pendientes que fija $\lambda$, de modo que el \emph{thrashing} es un impuesto constante y no una crisis.
  \item \textbf{Saber qué se sabe ayuda solo cuando evita pagar reparaciones, y no alcanza a no planificar.} \KB{} supera al lazo cerrado con el signo esperado en todo el rango, pero solo fuera del ruido con $r=3$ y $\lambda=0.5$ (\dif{+21.3}{3.8}); la ganancia desaparece en la ablación con horizonte \ahora, queda por debajo de \MB{} desde $\lambda=0.25$, y un solo tick por consulta invierte la ventaja. La monotonía, por su parte, no se rompe cuando el mundo se mueve, sino que se indexa por el tiempo.
  \item \textbf{Saber cuánto se sabe no mejoró la decisión.} \BY{} aprende $\lambda$ (error mediano de 0.020 con $r=3$) y su filtro está calibrado y tiene exactamente el soporte de la región lógica, pero queda por debajo del agente lógico (\dif{-16.0}{4.3} con $r=3$ y $\lambda=0.5$), porque su regla de utilidad esperada repara de más, y 83.0 puntos por debajo de \MB. El puente de la lógica a la probabilidad funciona en la creencia y falla en la decisión, cuyo modelo de costos resultó ser el eslabón débil.
  \item \textbf{Reparar vale menos, no más, cuando los recursos escasean.} Donde replanificar es gratis, la ventaja del lazo cerrado crece con la abundancia (de \dif{+36.7}{4.0} a \dif{+62.0}{5.0} a $\lambda=0.5$), y donde cuesta, crece la de no planificar, lo cual se explica por el costo de oportunidad de cada tick deliberando cuando sobran objetivos.
  \item \textbf{Lo que hay que medir de un despliegue es $\lambda$, y se puede medir.} $\lambda$ es el factor que más mueve a toda política con plan (hasta el 62\,\% de la varianza) y se recupera por ABC desde la telemetría de un motor en frío, lo cual basta para recomendar si planificar o no, mientras que $p$ no es identificable con esa telemetría ni decisivo para la recomendación.
\end{enumerate}

Volviendo a la pregunta, un agente que reescribe su plan al ver que el mundo cambió supera con claridad a uno que mantiene su plan hasta que falla solo cuando reescribir es casi gratis, el replanificador es bueno y el agente ve lo suficiente (con $c=0$, $p=1$ y visión total gana por entre 14 y 56 puntos según $\lambda$), mientras que con un costo de unos pocos ticks y visión parcial ambos casi empatan, y lo que sí vale desde $\lambda=0.1$ es conservar el trabajo en curso al reescribir, es decir, reparar en caliente en vez de reiniciar en frío. La ventaja de deliberar se invierte a partir de un costo de reescritura de unos pocos ticks (desde $c=2$ con $\lambda=0.5$, y entre $c=2$ y $c=5$ con $\lambda=0.25$, en unidades del modelo), pero no a favor del plan estático sino a favor de no planificar en absoluto, de forma que la decisión de cuándo reescribir resulta ser, antes que nada, la decisión de si deliberar, y se toma midiendo el dinamismo del entorno y el costo de deliberar, no la calidad del replanificador.
